\section{Model Merging: mejora de modelos sin entrenamiento}

\subsection{Conceptos básicos}

Model Merging es una técnica que combina directamente los parámetros de varios modelos, \textbf{sin necesidad de entrenamiento adicional}. Al principio se consideraba una «broma», pero ahora la ha adoptado toda empresa de LLM.

Idea central: tomar varios modelos de la misma arquitectura que pasaron por distintos fine-tuning y combinar sus parámetros mediante operaciones matemáticas (como el promedio ponderado), para obtener un nuevo modelo que reúna las ventajas de todos.

\begin{warningbox}{Limitaciones del Model Merging}
\begin{itemize}
    \item Solo pueden combinarse modelos de \textbf{la misma arquitectura y el mismo tamaño} (por ejemplo, solo se puede combinar Llama 3 con Llama 3, no entre generaciones distintas)
    \item Se recomienda usar modelos de \textbf{precisión completa} para la combinación: la pérdida de precisión de los modelos cuantizados se propaga al resultado combinado
    \item Elegir los pesos y parámetros adecuados se parece más a la «alquimia» que a una ciencia rigurosa: requiere experimentación repetida
\end{itemize}
\end{warningbox}

\subsection{Principales técnicas de combinación}

\subsubsection{SLERP (Spherical Linear Interpolation)}

SLERP realiza la interpolación sobre una esfera (en lugar de una interpolación lineal), lo que normalmente produce mejores resultados de combinación que el simple promedio lineal.

$$\text{SLERP}(\mathbf{p}_0, \mathbf{p}_1; t) = \frac{\sin((1-t)\Omega)}{\sin\Omega}\mathbf{p}_0 + \frac{\sin(t\Omega)}{\sin\Omega}\mathbf{p}_1$$

donde $\Omega = \arccos(\mathbf{p}_0 \cdot \mathbf{p}_1 / \|\mathbf{p}_0\| \|\mathbf{p}_1\|)$, $t \in {[}0, 1{]}$ es el coeficiente de interpolación.

Características: sencillo y confiable, muy popular en la comunidad de código abierto; pero \textbf{solo puede combinar dos modelos}.

\subsubsection{DARE + TIES}

Una técnica de combinación más avanzada, capaz de procesar más de dos modelos:

\begin{itemize}
    \item \textbf{DARE (Drop And REscale)}: poda aleatoriamente y reescala los parámetros de los modelos de origen
    \item \textbf{TIES (TrIm, Elect Sign, merge)}: conserva los parámetros más significativos y resuelve el problema de la cancelación de signos entre parámetros positivos y negativos mediante el \textbf{sign consensus} (consenso de signos)
\end{itemize}

Motivación del sign consensus: si un parámetro del modelo A es positivo y el mismo parámetro del modelo B es negativo, el simple promedio da un resultado cercano a cero, lo que en realidad hace perder la información que cada modelo aprendió por separado. TIES determina el signo del parámetro mediante un mecanismo de votación, evitando esta cancelación.

\subsection{Configuración práctica de Model Merging}

Con herramientas como \textbf{mergekit}, el esquema de combinación se define mediante un archivo de configuración YAML:

\begin{itemize}
    \item \textbf{density}: controla la proporción de parámetros que se conservan (en el método TIES no se conservan todos los parámetros)
    \item \textbf{weight}: controla el peso de cada modelo de origen en la combinación (se puede asignar un peso distinto según la puntuación de evaluación)
\end{itemize}

El expositor mostró el árbol genealógico del modelo Daredevil 8B: mediante combinación iterativa (un modelo combinado se vuelve a combinar con otros modelos) se puede construir un linaje de modelos complejo, y el producto final supera a todos los modelos de origen en varios benchmarks.

\subsection{Caso de aplicación: modelos específicos de idioma}

Proceso completo tomando como ejemplo la construcción de un LLM en finés:

\begin{enumerate}
    \item Preentrenamiento continuo sobre texto original en finés (5-100B tokens)
    \item SFT sobre instruction data en finés
    \item Preference Alignment sobre preference data en finés
    \item Problema: la capacidad del modelo en finés mejora mucho, pero \textbf{todas las demás capacidades se degradan gravemente}
    \item Solución: combinar el modelo en finés con un modelo instruct general (como Llama 3 Instruct)
    \item Resultado: se obtiene un modelo que \textbf{domina el finés y conserva las capacidades generales}
\end{enumerate}

\begin{importantbox}{El valor central de Model Merging}
Model Merging permite \textbf{superponer} las especialidades de distintos modelos sin sacrificar las capacidades que cada uno ya tenía. Es una estrategia eficaz para resolver el problema del «olvido catastrófico» (catastrophic forgetting) durante el fine-tuning, y resulta especialmente adecuada para escenarios como modelos específicos de idioma o modelos específicos de dominio.
\end{importantbox}

\subsection{Resumen de la sección}

Model Merging es una técnica de mejora de modelos sin costo de entrenamiento, que combina las ventajas de varios modelos mediante operaciones matemáticas a nivel de parámetros. SLERP es adecuado para combinar dos modelos, mientras que DARE+TIES admite la combinación de varios modelos. Es una herramienta poderosa para resolver el problema de degradación de capacidades causado por el fine-tuning.

\newpage
